% ============================================================
% Appendix Z.10 — FRFP Case Study (Clinical / Public Health):
% Autonomous AI for Diabetic Retinopathy Screening
% ============================================================

\section{FRFP Case Study (Clinical / Public Health):
  Autonomous AI for Diabetic Retinopathy Screening}
\label{app:frfp-dr-ai}

This appendix analyses autonomous AI systems for diabetic retinopathy (DR)
screening as a positive FRFP case study.
We focus primarily on FDA-cleared systems such as
LumineticsCore\textsuperscript{\textregistered} (formerly IDx-DR),
EyeArt\textsuperscript{\textregistered}, and
AEYE-DS\textsuperscript{\textregistered}.
These systems are among the first examples of autonomous AI diagnostics
deployed at scale in routine care, particularly in primary care clinics and
underserved settings.

\subsection{Background}

Diabetic retinopathy is a leading cause of preventable blindness among
working-age adults worldwide.
Regular retinal screening and timely treatment can prevent most vision loss,
but screening rates are chronically low, especially in primary care and
underserved populations.

To address this gap, several autonomous AI systems have been developed and
cleared:

\begin{itemize}
  \item \textbf{LumineticsCore\textsuperscript{\textregistered}
        (IDx-DR)}:
        in 2018, the first autonomous AI diagnostic system authorized by the
        U.S.\ FDA in any field of medicine.
        It automatically detects more-than-mild DR (mtmDR) from retinal
        images taken on a specific fundus camera in primary-care settings,
        without requiring an eye-care specialist to interpret each image.
  \item \textbf{EyeArt\textsuperscript{\textregistered}} (Eyenuk):
        FDA-cleared in 2020 as an autonomous AI system that detects both
        mtmDR and vision-threatening DR (vtDR) and can work across multiple
        fundus camera models.
  \item \textbf{AEYE-DS\textsuperscript{\textregistered}} (AEYE Health):
        FDA-cleared as an autonomous AI for DR screening, designed to
        operate with certain non-mydriatic fundus cameras and, in newer
        configurations, with a single image per eye and no dilation,
        targeting sub-minute screening workflows.
\end{itemize}

Pivotal trials for IDx-DR showed sensitivity around 87\% and specificity
around 90\% for mtmDR in primary care populations, meeting pre-specified
performance endpoints for autonomous use.
Subsequent implementation studies (e.g.\ at Johns Hopkins and in federally
qualified health centers) indicate that autonomous AI screening can
significantly increase screening rates and care-gap closure when integrated
into primary care workflows.

From an FRFP perspective, these systems are interesting because:

\begin{itemize}
  \item they \emph{are} autonomous diagnostic systems (they issue a
        clinical-level decision: refer vs.\ no-refer);
  \item yet they are deployed in ways that keep ultimate correctness at
        human and institutional boundaries (treatment, follow-up, programme
        design);
  \item their success depends heavily on careful explicit problem design,
        uncertainty labelling, and regulatory/institutional governance.
\end{itemize}

\subsection{Explicit and Tacit Spaces in DR Screening AI}

\subsubsection*{Explicit space \(E_{\mathrm{DR}}\)}

For systems such as LumineticsCore, EyeArt, and AEYE-DS, explicit space
includes:

\begin{itemize}
  \item input images:
        \begin{itemize}
          \item colour fundus photographs from specified camera models;
          \item typically macula- and disc-centered views per eye;
        \end{itemize}
  \item model internals:
        \begin{itemize}
          \item trained deep convolutional or related networks;
          \item pre-processing and image-quality assessment modules.
        \end{itemize}
  \item output artefacts:
        \begin{itemize}
          \item a categorical diagnostic conclusion (e.g.\ ``mtmDR
                present'' vs.\ ``mtmDR absent'');
          \item image-quality verdict (gradable vs.\ ungradable);
          \item sometimes a DR severity grade or referral recommendation.
        \end{itemize}
  \item deployment infrastructure:
        \begin{itemize}
          \item cloud or local servers receiving images and returning
                decisions;
          \item user interfaces in primary care clinics and retail
                settings.
        \end{itemize}
\end{itemize}

All these artefacts live in \(E_{\mathrm{DR}}\): they are explicit bits and
decisions produced and consumed by machines.

\subsubsection*{Tacit space \(T_{\mathrm{DR}}\)}

Tacit space contains the epistemic states of:

\begin{itemize}
  \item primary-care clinicians:
        their understanding of diabetes care pathways, risk factors, and
        when to emphasize follow-up;
  \item eye-care specialists (ophthalmologists, retina specialists):
        nuanced knowledge of DR progression, treatment thresholds, and
        patient-specific considerations;
  \item patients:
        their understanding of vision risk, willingness to attend
        referrals, and trust in the healthcare system;
  \item health-system administrators and payers:
        beliefs about programme design, reimbursement, and equity.
\end{itemize}

These tacit states determine how explicit AI outputs are interpreted and
acted upon: whether a ``refer'' result leads to timely specialist care, or
whether a ``no-refer'' result alters behaviour around future screening.

\subsection{Boundary Design:
  From Image Classification to Clinical Workflow}

\subsubsection*{Regulatory definition of the AI boundary}

Regulatory documents for IDx-DR and similar systems define a very specific
boundary:

\begin{itemize}
  \item The AI system is authorised to \emph{autonomously} detect more-than-
        mild DR (and for EyeArt/AEYE-DS, in some cases, vtDR) from images
        acquired under specified conditions.
  \item If image quality is insufficient, the system returns an
        ``ungradable'' result and the patient must be referred or re-imaged.
  \item For negative results (no mtmDR detected), guidelines and labelling
        still require repeat screening at defined intervals (often yearly),
        not discharge from care.
  \item For positive results, the system issues a referral recommendation,
        but an ophthalmologist ultimately determines diagnosis and
        treatment.
\end{itemize}

In FRFP terms:

\begin{itemize}
  \item the AI’s explicit decision boundary is: classify image $\to$
        \{mtmDR present, mtmDR absent, ungradable\};
  \item the clinical boundary---where correctness is endorsed---remains:
        \begin{itemize}
          \item in primary care: deciding whether a patient has satisfied
                the screening requirement or needs referral;
          \item in specialist care: diagnosing DR stage and choosing
                treatment (laser, anti-VEGF, observation, etc.).
        \end{itemize}
\end{itemize}

The FDA explicitly recognises this structure by authorising the AI to make a
narrow, well-defined diagnostic decision while leaving broader care
decisions to human clinicians.

\subsubsection*{Workflow integration as boundary implementation}

Implementation studies describe how DR AI is embedded in primary-care
workflows:

\begin{itemize}
  \item a nurse or medical assistant captures fundus images during a routine
        diabetes visit;
  \item the AI system returns a result in minutes;
  \item if positive or ungradable, staff schedule an eye-care referral and
        provide counselling;
  \item if negative, the result is documented in the EHR, closing the DR
        screening care gap for that year.
\end{itemize}

From an FRFP perspective:

\begin{itemize}
  \item AI operates fully within explicit space (image $\to$ label), but
        boundary events (referral decisions, treatment) are explicitly
        structured in clinical workflows and guidelines;
  \item the architecture intentionally shifts some correctness
        responsibility (for \emph{screening} classification) into the AI
        pipeline, but within a tightly circumscribed domain and with
        oversight via quality metrics and regulatory conditions.
\end{itemize}

\subsection{Protocol-Level Properties in FRFP Terms}

\subsubsection*{Z.10.1 Explicit Semantics:
  A Narrow, Well-Specified Task}

DR screening AI enjoys strong explicit semantics:

\begin{itemize}
  \item \textbf{Task definition}:
        detect the presence of more-than-mild DR (and sometimes vtDR) from
        retinal photographs, according to accepted grading scales.
  \item \textbf{Gold standard}:
        indendent, masked grading of images by retina specialists serves as
        a clear ground truth in pivotal trials.
  \item \textbf{Performance targets}:
        pre-specified sensitivity/specificity thresholds for autonomous use
        (e.g.\ rule-out sensitivity, rule-in specificity) are agreed with
        regulators.
\end{itemize}

In FRFP language:

\begin{itemize}
  \item the mapping
        \(
          f_{\mathrm{DR}} : \text{fundus image} \to
          \{\text{mtmDR}, \text{no mtmDR}, \text{ungradable}\}
        \)
        has a precise semantics under trial conditions;
  \item $E_{\mathrm{DR}}$ has reliable, externally validated metrics
        to assess correctness of the AI’s explicit outputs;
  \item this is exactly the kind of narrow medical task where explicit AI
        can legitimately assume part of the correctness burden.
\end{itemize}

\subsubsection*{Z.10.2 Uncertainty and ``Ungradable'' as Structural Safeguards}

An important design feature is the explicit handling of uncertainty:

\begin{itemize}
  \item systems classify some images as ungradable, prompting referral to an
        eye-care provider instead of issuing a false negative;
  \item image-quality checks and fail-safe logic are part of the cleared
        indication for use;
  \item in practice, this means that when the explicit pipeline cannot
        produce a reliable decision, it falls back to tacit human judgement.
\end{itemize}

In FRFP terms:

\begin{itemize}
  \item ``ungradable'' is a designed escape hatch back into tacit space
        \(T_{\mathrm{DR}}\);
  \item the explicit system is not forced to speak when its semantics would
        be unreliable;
  \item this is a concrete instantiation of FRFP’s recommendation that
        explicit pipelines carry their own limitations to the boundary.
\end{itemize}

\subsubsection*{Z.10.3 Safe Horizons in Screening Programs}

DR screening is a repeated, population-level process:

\begin{itemize}
  \item patients with diabetes should be screened annually or every
        few years, depending on guidelines;
  \item autonomous AI is intended for long-horizon use across large
        populations (health systems, national programmes).
\end{itemize}

FRFP’s safe-horizon lens highlights several structural properties:

\begin{itemize}
  \item \textbf{Temporal horizons}:
        even after a negative AI screen, guidelines insist on re-screening
        after a fixed interval; no sequence of negative outputs permanently
        discharges the patient from surveillance.
  \item \textbf{Data shift monitoring}:
        camera changes, population differences, and disease prevalence
        shifts can degrade performance; responsible deployments monitor
        quality metrics and periodically revalidate the AI.
  \item \textbf{Programme-level review}:
        implementation studies treat DR AI roll-outs as quality-improvement
        projects, with ongoing measurement of screening uptake,
        time-to-referral, and adherence.
\end{itemize}

Thus, even though the AI is autonomous at the single-visit level, the
overall programme is embedded in human and institutional cycles that limit
how far the system can run without tacit correction.

\subsection{Institutional Semantics:
  Regulation, Reimbursement, and Equity}

\subsubsection*{Regulatory governance as FRFP operator}

Regulators act as institutional operators in FRFP terms:

\begin{itemize}
  \item The FDA’s De Novo authorization of IDx-DR set a precedent for
        autonomous AI in medicine, including requirements on performance,
        indications, and image-acquisition conditions.
  \item Subsequent 510(k) clearances for EyeArt and AEYE-DS followed,
        broadening the set of devices and cameras but within a shared
        regulatory framework.
  \item CPT codes and reimbursement policies have been developed to support
        billing for autonomous AI DR screening, reinforcing its role in
        routine care.
\end{itemize}

These institutional choices:

\begin{itemize}
  \item formalise the boundary between AI output and clinical use;
  \item define safe deployment domains (primary care, camera types,
        age ranges);
  \item require ongoing post-market surveillance and device updates.
\end{itemize}

\subsubsection*{Health-system deployments and equity goals}

Large health systems and public programmes increasingly deploy DR AI with
explicit equity aims:

\begin{itemize}
  \item Johns Hopkins Medicine’s deployment across primary-care sites showed
        increased closure of screening care gaps, especially in some
        underserved groups.
  \item Multicomponent trials in federally qualified health centres
        (FQHCs) evaluate whether AI screening plus EHR integration and
        patient education improves adherence to screening and standards of
        diabetes care.
  \item In India, AI-based DR screening platforms such as MadhuNetrAI /
        MadhuNETrAI are being piloted or rolled out in government
        programmes to reach rural and military populations with limited
        ophthalmology access.
\end{itemize}

From an FRFP perspective, these programmes:

\begin{itemize}
  \item treat autonomous AI as part of an institutional operator that
        adjusts population-level tacit states (e.g.\ awareness, screening
        norms, referral patterns);
  \item rely on explicit AI outputs to trigger care, but leave diagnosis,
        treatment, and policy adaptation in human hands;
  \item illustrate how explicit systems can be embedded in equity-oriented
        institutional dynamics, rather than only in revenue-maximising
        workflows.
\end{itemize}

\subsection{Counterfactual Pitfalls and FRFP Safeguards}

DR screening AI also highlights potential FRFP failure modes that are being
actively managed:

\begin{itemize}
  \item \textbf{Overextension of indications}.
        Using an AI system beyond its cleared population (e.g.\ different
        ethnic mix, camera, or disease spectrum) without revalidation could
        degrade performance.
        FRFP would insist on explicit safe horizons and re-grounding when
        context shifts.
  \item \textbf{Silent false negatives}.
        If negative AI results led clinicians or patients to skip future
        exams, long-horizon risk would increase.
        Current guidelines and labelling prevent this by mandating
        periodic re-screening regardless of AI output.
  \item \textbf{Complacency in specialist care}.
        There is a risk that specialists might over-trust AI-based triage
        and pay less attention to borderline or atypical cases.
        FRFP would recommend clear communication that AI covers only the
        screening threshold, not full disease management.
\end{itemize}

So far, regulatory frameworks, trial designs, and deployment practices have
kept many of these issues under control, making DR screening AI a relatively
successful instance of autonomy within an FRFP-compatible architecture.

\subsection{Lessons for FRFP and Clinical AI}

Autonomous DR screening systems illustrate several of FRFP’s central points
in a positive light:

\begin{enumerate}
  \item \textbf{Narrow, well-specified tasks can bear autonomous decisions.}
        When explicit semantics are strong (image $\to$ DR threshold
        classification) and ground truth is available, AI can legitimately
        take on a clearly-bounded diagnostic role.

  \item \textbf{Explicit uncertainty and ``ungradable'' outputs are crucial.}
        Structured fallbacks to human care are built into the model’s
        explicit interface, embodying FRFP’s recommendation that pipelines
        carry their own limits back to the boundary.

  \item \textbf{Safe horizons can be encoded as clinical intervals.}
        Annual (or guideline-based) re-screening and ongoing performance
        monitoring function as safe-horizon mechanisms in a public-health
        setting.

  \item \textbf{Institutional governance makes autonomy possible.}
        Regulatory authorisation, reimbursement structures, and programme
        design act as FRFP-style institutional operators that shape how AI
        is used and prevent over-automation of tacit clinical judgment.

  \item \textbf{Explicit AI can advance equity when embedded correctly.}
        By bringing DR screening to primary care, FQHCs, and low-resource
        settings, DR AI demonstrates how explicit systems can extend access
        while keeping correctness at human and institutional boundaries.
\end{enumerate}

As such, autonomous DR screening is a valuable counterpoint to high-profile
failures:
it shows that, under FRFP-compatible conditions---narrow task, strong
semantics, explicit uncertainty, and robust institutional governance---AI
can safely and measurably improve population health outcomes.
